% !TeX program = xelatex
% Informe de Proyecto - FisioCare
% Plantilla adaptada del proyecto de Overleaf del curso

\documentclass[12pt,a4paper,oneside]{report}

\usepackage[spanish,es-nodecimaldot]{babel}
\usepackage{fontspec}
\setmainfont{Noto Serif}
\setsansfont{Lato}
\setmonofont{DejaVu Sans Mono}
\usepackage[a4paper,top=2.5cm,bottom=2.4cm,left=2.7cm,right=2.5cm,headheight=16pt]{geometry}
\usepackage{setspace}
\onehalfspacing
\usepackage{microtype}
\usepackage{graphicx}
\usepackage{float}
\usepackage{booktabs}
\usepackage{array}
\usepackage{tabularx}
\usepackage{longtable}
\usepackage{multirow}
\usepackage{enumitem}
\usepackage{amsmath}
\usepackage{amssymb}
\usepackage[table]{xcolor}
\usepackage{fancyhdr}
\usepackage{titlesec}
\usepackage{tocloft}
\usepackage{caption}
\usepackage{tikz}
\usetikzlibrary{positioning,arrows.meta,shapes.geometric,fit,calc}
\usepackage{hyperref}
\usepackage{bookmark}
\usepackage{url}
\usepackage{ragged2e}
\usepackage[most]{tcolorbox}

\definecolor{CibertecBlue}{HTML}{123B63}
\definecolor{CibertecLight}{HTML}{EAF1F7}
\definecolor{CibertecAccent}{HTML}{2E74B5}
\definecolor{TextGray}{HTML}{4D4D4D}
\definecolor{LightGray}{HTML}{F3F5F7}
\definecolor{LineGray}{HTML}{C9D2DA}
\definecolor{CibertecInk}{HTML}{23384D}

\hypersetup{
  colorlinks=true,
  linkcolor=CibertecBlue,
  urlcolor=CibertecAccent,
  citecolor=CibertecBlue,
  pdftitle={FisioCare - Plataforma web para reservas de fisioterapia},
  pdfauthor={Equipo FisioCare},
  pdfsubject={Proyecto de Investigación Aplicada}
}
\urlstyle{same}

\setlength{\parindent}{1.25cm}
\setlength{\parskip}{0pt}
\setlength{\emergencystretch}{2em}
\setlength{\tabcolsep}{6pt}
\renewcommand{\arraystretch}{1.22}
\setlist[itemize]{leftmargin=1.1cm,itemsep=2pt,topsep=4pt}
\setlist[enumerate]{leftmargin=1.1cm,itemsep=2pt,topsep=4pt}
\renewcommand{\labelitemi}{\textcolor{CibertecAccent}{\small\textbullet}}

\pagestyle{fancy}
\fancyhf{}
\fancyhead[L]{\small\sffamily\color{CibertecInk}Informe de Proyecto - FisioCare}
\fancyhead[R]{\small\sffamily\bfseries\color{CibertecBlue}CIBERTEC}
\fancyfoot[L]{\small\sffamily\color{TextGray}Computación e Informática}
\fancyfoot[R]{\small\sffamily\color{TextGray}Página \thepage}
\renewcommand{\headrulewidth}{0.5pt}
\renewcommand{\footrulewidth}{0.5pt}
\fancypagestyle{plain}{
  \fancyhf{}
  \fancyhead[L]{\small\sffamily\color{CibertecInk}Informe de Proyecto - FisioCare}
  \fancyhead[R]{\small\sffamily\bfseries\color{CibertecBlue}CIBERTEC}
  \fancyfoot[L]{\small\sffamily\color{TextGray}Computación e Informática}
  \fancyfoot[R]{\small\sffamily\color{TextGray}Página \thepage}
  \renewcommand{\headrulewidth}{0.5pt}
  \renewcommand{\footrulewidth}{0.5pt}
}

\titleformat{\chapter}[display]
  {\normalfont\sffamily\bfseries\color{CibertecBlue}}
  {\fontsize{14}{16}\selectfont CAPÍTULO \thechapter}
  {6pt}
  {\fontsize{18}{22}\selectfont\MakeUppercase}
  [\vspace{6pt}\color{CibertecAccent}\titlerule]
\titlespacing*{\chapter}{0pt}{-18pt}{24pt}
\titleformat{\section}{\normalfont\sffamily\bfseries\large\color{CibertecBlue}}{\thesection.}{0.6em}{}[\vspace{2pt}{\color{CibertecLight}\titlerule}]
\titleformat{\subsection}{\normalfont\sffamily\bfseries\normalsize\color{CibertecAccent}}{\thesubsection.}{0.6em}{}
\titleformat{\subsubsection}{\normalfont\sffamily\bfseries\normalsize\color{TextGray}}{}{0em}{}
\renewcommand{\contentsname}{ÍNDICE}
\setcounter{tocdepth}{2}
\setcounter{secnumdepth}{3}
\renewcommand{\cftchapfont}{\bfseries\sffamily\color{CibertecBlue}}
\renewcommand{\cftchappagefont}{\bfseries\color{CibertecBlue}}
\renewcommand{\cftsecfont}{\sffamily}
\renewcommand{\cftsubsecfont}{\sffamily\small}
\setlength{\cftbeforechapskip}{7pt}

\newcolumntype{Y}{>{\raggedright\arraybackslash}X}
\newcommand{\source}[1]{\par\noindent{\footnotesize\textcolor{TextGray}{#1}}\par}
\newcommand{\ciberteclogo}{cibertec-logo-png_seeklogo-315758.png}
\newtcolorbox{projectsummary}{
  enhanced,breakable,
  colback=CibertecLight,colframe=CibertecAccent,
  boxrule=0.7pt,arc=4pt,
  left=12pt,right=12pt,top=10pt,bottom=10pt,
  title={Resumen ejecutivo},
  fonttitle=\sffamily\bfseries\color{CibertecBlue},
  coltitle=CibertecBlue
}

\begin{document}

\begin{titlepage}
\thispagestyle{empty}
\begin{tikzpicture}[remember picture,overlay]
  \fill[CibertecBlue] (current page.north west) rectangle ([xshift=2.8cm]current page.south west);
  \fill[CibertecAccent] ([xshift=2.8cm]current page.north west) rectangle ([xshift=3.05cm]current page.south west);
  \fill[CibertecLight] ([xshift=3.05cm,yshift=-1.3cm]current page.north west)
    rectangle (current page.north east);
  \fill[CibertecAccent] ([xshift=3.05cm,yshift=-1.3cm]current page.north west)
    rectangle ([xshift=3.05cm,yshift=-1.38cm]current page.north east);
\end{tikzpicture}

\hspace*{1.1cm}
\begin{minipage}{0.82\textwidth}
\centering
\vspace*{0.55cm}
\IfFileExists{\ciberteclogo}{\includegraphics[width=5.4cm]{\ciberteclogo}\par}{\sffamily\bfseries\Large\color{CibertecBlue}CIBERTEC\par}
\vspace{0.65cm}
{\sffamily\small\color{CibertecAccent}\MakeUppercase{Carrera profesional: Computación e Informática}\par}
\vspace{1.25cm}
{\sffamily\bfseries\color{CibertecBlue}\fontsize{32}{38}\selectfont FISIOCARE\par}
\vspace{0.35cm}
{\sffamily\color{TextGray}\fontsize{15}{20}\selectfont Plataforma web para la gestión\par}
{\sffamily\color{TextGray}\fontsize{15}{20}\selectfont de reservas de fisioterapia\par}
\vspace{0.8cm}
\begin{tcolorbox}[width=0.92\textwidth,colback=CibertecBlue,colframe=CibertecBlue,boxrule=0pt,arc=3pt,left=10pt,right=10pt,top=8pt,bottom=8pt]
\centering\sffamily\bfseries\color{white}\large Proyecto de Investigación Aplicada
\end{tcolorbox}
\vspace{0.7cm}
{\sffamily\bfseries\color{CibertecBlue}\large Desarrollo de Aplicaciones Web I (4694)\par}
\vspace{0.15cm}
{\sffamily\color{TextGray}\normalsize Quinto ciclo \textbullet\ Semestre 2026-I\par}
\vfill
\begin{tcolorbox}[width=0.96\textwidth,colback=LightGray,colframe=LineGray,boxrule=0.6pt,arc=3pt,left=12pt,right=12pt,top=10pt,bottom=10pt]
\begin{tabularx}{\textwidth}{@{}p{3.1cm}Y@{}}
\textcolor{CibertecBlue}{\sffamily\bfseries Docente} & Berrio Huamani, Miguel Angel\\[7pt]
\textcolor{CibertecBlue}{\sffamily\bfseries Alumnos} & Alcarraz Ortiz, Angel Máximo\\
& Canazas Rosas, Willian Roy
\end{tabularx}
\end{tcolorbox}
\vfill
{\sffamily\color{TextGray}\small Lima, Perú \textbullet\ 2026\par}
\end{minipage}
\end{titlepage}

\pagenumbering{roman}
\tableofcontents
\clearpage
\pagenumbering{arabic}

\chapter{Proyecto de Investigación}

\section{Resumen}
\begin{projectsummary}
FisioCare es una aplicación web empresarial para organizar la reserva y el seguimiento de sesiones de fisioterapia. Permite que el paciente conozca los servicios, cree una cuenta, inicie sesión con correo o Google, registre una reserva y revise su estado. El administrador gestiona la agenda, confirma o rechaza solicitudes, cancela o elimina registros y administra el catálogo.

El backend se construyó con Java y Spring Boot, Spring MVC, Spring Data JPA y Spring Security. El frontend se desarrolló con Angular y consume una API REST. PostgreSQL persiste usuarios, servicios y reservas. Las contraseñas locales se protegen mediante hashing BCrypt y el acceso se gestiona con JWT. Docker Compose permite ejecutar la solución completa.
\end{projectsummary}

\section{Introducción}
La coordinación de una atención fisioterapéutica puede involucrar mensajes dispersos, confirmaciones manuales y poca visibilidad del estado de cada solicitud. FisioCare centraliza el proceso en una plataforma web responsive con flujos diferenciados para pacientes y administradores.

El proyecto busca facilitar el acceso a los servicios de rehabilitación y reducir el trabajo manual asociado a la agenda. El impacto esperado es una atención más ordenada, trazable y accesible para una clínica o consultorio de fisioterapia.

\section{Diagnóstico: análisis SEPTE}
El análisis SEPTE prioriza las variables Social, Tecnológica y Política/Legal, porque son las que se relacionan directamente con la oportunidad de mejora.

\begin{longtable}{|p{2.4cm}|p{7.0cm}|p{6.0cm}|}
\caption{Análisis SEPTE aplicado a FisioCare}\label{tab:septe}\\
\hline
\rowcolor{CibertecBlue}\color{white}\textbf{Variable} & \color{white}\textbf{Hallazgo y evidencia} & \color{white}\textbf{Implicación para FisioCare}\\
\hline
\endfirsthead
\hline
\rowcolor{CibertecBlue}\color{white}\textbf{Variable} & \color{white}\textbf{Hallazgo y evidencia} & \color{white}\textbf{Implicación para FisioCare}\\
\hline
\endhead
\rowcolor{CibertecLight}
Social & La OMS estima que 2,4 mil millones de personas viven con una condición que podría beneficiarse de rehabilitación.\footnote{OMS, ``Rehabilitación'', 2024. \url{https://www.who.int/es/news-room/fact-sheets/detail/rehabilitation}} & Facilitar la solicitud de cita, mostrar estados y orientar el siguiente paso. Riesgo: brecha de acceso, baja alfabetización digital o abandono.\\ \hline
Tecnológica & El INEI reportó acceso a Internet en hogares del país y una diferencia importante entre Lima Metropolitana y el área rural.\footnote{INEI, ``Las Tecnologías de Información y Comunicación en los Hogares'', II trimestre 2024. \url{https://www.inei.gob.pe/media/MenuRecursivo/boletines/boletin_tic_iit24.pdf}} & Interfaz responsive, ligera y usable desde celular. Riesgo: conectividad desigual y dependencia del backend.\\ \hline
Política / legal & La Ley N.° 29733 protege el derecho fundamental a la protección de datos personales en Perú.\footnote{Congreso de la República del Perú, Ley N.° 29733. \url{https://www.gob.pe/institucion/congreso-de-la-republica/normas-legales/243470-29733}} & Autenticación, autorización por roles y hash de contraseñas. Riesgo: exposición de datos si se despliega sin HTTPS y controles.\\ \hline
\end{longtable}

\section{Objetivos}
\subsection{Objetivo general}
Desarrollar una aplicación web empresarial que centralice la reserva y gestión de sesiones de fisioterapia mediante un backend Java con Spring Framework, una API REST persistente y un frontend Angular.

\subsection{Objetivos específicos SMART}
\begin{enumerate}
\item Implementar antes de la sustentación final una API REST protegida que permita autenticar usuarios, gestionar servicios y ejecutar GET, POST, PUT y DELETE sobre recursos persistidos en PostgreSQL.
\item Validar antes de la entrega final las operaciones de insertar, actualizar, eliminar y listar mediante pruebas automatizadas de la capa de acceso a datos.
\item Permitir que un paciente registre una reserva y consulte su estado, y que un administrador la confirme, rechace, cancele o elimine desde el panel de gestión.
\end{enumerate}

\section{Justificación del proyecto}
FisioCare transforma un proceso sensible a errores en un flujo digital trazable. La reserva conserva paciente, servicio, fecha, estado y responsable de la gestión. La separación entre Angular, Spring Boot, JPA y PostgreSQL permite mantener y probar el sistema.

\subsection{Beneficiarios directos}
\begin{itemize}
\item Pacientes que desean consultar servicios y registrar una cita desde una landing page.
\item Personal administrativo o fisioterapeutas que gestionan la agenda en un panel.
\item Equipo desarrollador que aplica Spring, Angular, REST, seguridad, pruebas y Docker.
\end{itemize}

\subsection{Beneficiarios indirectos}
\begin{itemize}
\item Familiares o acompañantes que apoyan la coordinación de citas.
\item La clínica o consultorio, al disponer de información ordenada para mejorar su atención.
\end{itemize}

\section{Definición y alcance}
El sistema cubre el flujo desde la presentación del servicio hasta la gestión de una reserva. Incluye registro e inicio de sesión local, login opcional con Google OAuth2, autorización por roles CUSTOMER y ADMIN, catálogo de servicios, disponibilidad básica, creación y consulta de reservas, confirmación, rechazo, cancelación y eliminación.

No incluye historia clínica, diagnóstico médico, pagos en línea, videollamadas, notificaciones por correo ni facturación. Estas funciones se consideran extensiones futuras.

\subsection{Arquitectura y tecnologías}
\rowcolors{2}{CibertecLight}{white}
\begin{tabularx}{\textwidth}{|p{3cm}|p{4.2cm}|Y|}
\hline
\rowcolor{CibertecBlue}\color{white}\textbf{Capa} & \color{white}\textbf{Tecnología} & \color{white}\textbf{Responsabilidad}\\
\hline
Frontend & Angular & Landing, autenticación, reserva, cuenta del paciente y panel administrador.\\ \hline
Backend & Java 17 + Spring Boot & Controladores REST, validaciones, reglas, seguridad y JWT.\\ \hline
Persistencia & Spring Data JPA + PostgreSQL & Usuarios, servicios y reservas.\\ \hline
Seguridad & Spring Security + BCrypt + OAuth2/OIDC & Roles, hash de contraseñas, JWT y Google.\\ \hline
Ejecución & Docker Compose & Contenedores para frontend, backend y base de datos.\\ \hline
\end{tabularx}
\rowcolors{1}{white}{white}

\subsection{Endpoints REST principales}
\rowcolors{2}{CibertecLight}{white}
\begin{tabularx}{\textwidth}{|p{2cm}|p{5.2cm}|Y|}
\hline
\rowcolor{CibertecBlue}\color{white}\textbf{Método} & \color{white}\textbf{Endpoint} & \color{white}\textbf{Uso}\\
\hline
POST & /api/auth/register & Registrar usuario con contraseña protegida mediante BCrypt.\\ \hline
POST & /api/auth/login & Autenticar y emitir JWT.\\ \hline
GET & /api/auth/me & Consultar el usuario autenticado.\\ \hline
GET & /api/auth/providers & Consultar proveedores de autenticación disponibles.\\ \hline
GET & /api/services & Listar servicios disponibles.\\ \hline
POST & /api/services & Crear servicio como administrador.\\ \hline
PUT & /api/services/\{id\} & Actualizar servicio.\\ \hline
DELETE & /api/services/\{id\} & Eliminar servicio.\\ \hline
GET & /api/reservations & Listar reservas como administrador.\\ \hline
GET & /api/reservations/mine & Listar reservas del paciente autenticado.\\ \hline
GET & /api/reservations/availability & Consultar disponibilidad por servicio y fecha.\\ \hline
POST & /api/reservations & Crear reserva.\\ \hline
PUT & /api/reservations/\{id\} & Reprogramar reserva.\\ \hline
PATCH & /api/reservations/\{id\}/status & Confirmar o rechazar una reserva.\\ \hline
PATCH & /api/reservations/\{id\}/cancel & Cancelar una reserva.\\ \hline
DELETE & /api/reservations/\{id\} & Eliminar reserva.\\ \hline
\end{tabularx}
\rowcolors{1}{white}{white}

\section{Productos y entregables}
\begin{itemize}
\item Código fuente del backend, frontend y configuración Docker Compose.
\item Base de datos PostgreSQL persistida mediante volumen Docker.
\item Informe de proyecto en formato editable.
\item Presentación para la exposición.
\item Guion estructurado para video demo de 3 a 5 minutos.
\item Pruebas automatizadas de persistencia y evidencias de build.
\end{itemize}

\section{Conclusiones}
\begin{enumerate}
\item La solución integra Spring Boot y Angular mediante servicios REST persistentes, con operaciones GET, POST, PUT y DELETE.
\item La separación de roles y el uso de BCrypt y JWT fortalecen el control de acceso; OAuth2 agrega una alternativa conveniente.
\item La digitalización de la agenda mejora la trazabilidad y deja una base para incorporar notificaciones, pagos e historia clínica.
\end{enumerate}

\section{Recomendaciones}
\begin{enumerate}
\item Desplegar con HTTPS, secretos administrados por el proveedor cloud y una política de privacidad antes de usar datos reales.
\item Incorporar pruebas de integración REST y pruebas end-to-end del flujo paciente-administrador.
\item Agregar recordatorios, disponibilidad por profesional y calendario externo después de validar el proceso operativo.
\end{enumerate}

\clearpage
\section{Glosario}
\begin{tabularx}{\textwidth}{|>{\raggedright\arraybackslash\bfseries}p{3.7cm}|Y|}
\hline
\rowcolor{CibertecBlue}\color{white}\textbf{Término} & \color{white}\textbf{Definición}\\
\hline
API REST & Interfaz que expone recursos mediante HTTP y métodos GET, POST, PUT y DELETE.\\
\hline
Angular & Framework frontend basado en TypeScript.\\
\hline
BCrypt & Algoritmo de hashing adaptativo usado para proteger contraseñas.\\
\hline
Docker Compose & Herramienta que coordina los contenedores del frontend, backend y base de datos.\\
\hline
JWT & JSON Web Token firmado que transporta identidad y permisos.\\
\hline
OAuth2/OIDC & OAuth2 delega la autorización y OIDC añade identidad sobre ese flujo, como en el acceso con Google.\\
\hline
JPA & Especificación Java para mapear objetos a tablas.\\
\hline
PostgreSQL & Sistema gestor de base de datos relacional usado para persistir usuarios, servicios y reservas.\\
\hline
CRUD & Crear, leer, actualizar y eliminar información.\\
\hline
Reserva & Solicitud de atención asociada a un paciente, servicio, fecha y estado.\\
\hline
\end{tabularx}

\clearpage
\section{Bibliografía}
\begin{enumerate}[label={\textcolor{CibertecAccent}{[\arabic*]}},leftmargin=1.1cm,itemsep=7pt]
\item Organización Mundial de la Salud. (2024). \textit{Rehabilitación}. \url{https://www.who.int/es/news-room/fact-sheets/detail/rehabilitation}
\item Instituto Nacional de Estadística e Informática. (2024). \textit{Las Tecnologías de Información y Comunicación en los Hogares}. \url{https://www.inei.gob.pe/estadisticas/indice-tematico/tecnologias-de-la-informacion-y-telecomunicaciones/}
\item Congreso de la República del Perú. \textit{Ley N.° 29733, Ley de Protección de Datos Personales}. \url{https://www.gob.pe/institucion/congreso-de-la-republica/normas-legales/243470-29733}
\item Spring. \textit{Spring Boot Reference Documentation}. \url{https://docs.spring.io/spring-boot/reference/}
\item Angular. \textit{Angular Documentation}. \url{https://angular.dev/}
\end{enumerate}

\appendix
\chapter{Anexos}
\section{Evidencia técnica de cumplimiento}
\rowcolors{2}{CibertecLight}{white}
\begin{tabularx}{\textwidth}{|>{\raggedright\arraybackslash}p{3.8cm}|Y|>{\centering\arraybackslash}p{3.0cm}|}
\hline
\rowcolor{CibertecBlue}\color{white}\textbf{Criterio} & \color{white}\textbf{Evidencia en FisioCare} & \color{white}\textbf{Estado}\\
\hline
Login REST y BCrypt & AuthController, PasswordEncoder y usuarios en PostgreSQL. & Implementado\\
\hline
Operaciones REST & BookingServiceController y ReservationController. & Implementado\\
\hline
Consumo REST desde Angular & BookingApiService con GET, POST, PUT, PATCH y DELETE. & Implementado\\
\hline
Persistencia & Entidades JPA y repositorios. & Implementado\\
\hline
Pruebas de datos & ReservationRepositoryDataTest: insertar, listar, actualizar y eliminar. & Implementado\\
\hline
Contenedores & compose.yaml con frontend, backend y PostgreSQL. & Implementado\\
\hline
\end{tabularx}
\rowcolors{1}{white}{white}

\section{Credenciales de demostración local}
Administrador: \texttt{admin@reservas.local} / \texttt{Admin123!}.\\
Cliente demo: \texttt{cliente@reservas.local} / \texttt{Cliente123!}.\\
\textcolor{TextGray}{\small Estas credenciales son únicamente para la demostración local y no deben utilizarse en producción.}

\section{Guion resumido del video demo}
{\small
\begin{enumerate}[itemsep=1pt,topsep=2pt]
\item Presentar el problema y la propuesta FisioCare (30 segundos).
\item Mostrar landing, servicios y Reservar cita (40 segundos).
\item Registrar una reserva y mostrar Mi cuenta (60 segundos).
\item Gestionar la reserva y editar un servicio como administrador (80 segundos).
\item Mostrar Docker, API REST y pruebas Maven; cerrar con conclusiones (40 segundos).
\end{enumerate}
}

\end{document}
